% ======================================================================
% EXPANSÃO — Impacto do SDD+TDD no Ecossistema
% ======================================================================

\section{O Impacto Metodológico do SDD+TDD no OpenCode Ecosystem}

\subsection{Antes do SDD+TDD: Correção Reativa (Rounds 1--7)}

Nos primeiros sete rounds de desenvolvimento, o ecossistema OpenCode operava
sem uma metodologia formal de teste e verificação. O paradigma dominante era
a \textbf{correção reativa}: o sistema produzia um artefato (artigo acadêmico,
análise de dados, busca de editais), uma banca simulada de revisores
identificava problemas, e um loop de correção iterava manualmente até que o
resultado fosse considerado aceitável.

Esta abordagem, embora funcional — a progressão de scores de 85 para 98
atesta sua eficácia relativa —, apresentava três limitações estruturais que
se tornariam evidentes à medida que a complexidade do ecossistema aumentava:

\textbf{Primeiro, ausência de garantia contra regressões.} Cada correção
aplicada a um artigo ou pipeline podia, inadvertidamente, quebrar outras
partes do sistema. O Round 7 (editais-br v7.1) exemplificou esta fragilidade:
um bug \texttt{KeyError} no sistema de scoring de editais só foi descoberto
porque um operador humano inspecionou manualmente a saída. Não havia
\textit{nenhum} teste automatizado que detectaria a regressão.

\textbf{Segundo, rastreabilidade limitada.} As decisões de correção — por que
uma fórmula foi alterada, qual o racional para substituir um termo — residiam
na memória de curto prazo da sessão. Não havia registro estruturado que
permitisse a um agente (ou a um humano) entender, semanas depois, por que
uma escolha específica foi feita.

\textbf{Terceiro, dependência de supervisão humana.} O loop de correção
exigia que um operador validasse cada iteração. Para um sistema que aspirava à
autonomia, esta dependência representava um gargalo fundamental: a velocidade
de evolução do ecossistema era limitada pela velocidade de revisão humana.

A Tabela~\ref{tab:pre_sdd} resume as características do ecossistema antes da
introdução do SDD+TDD.

\begin{table}[H]\centering
\caption{Ecossistema OpenCode antes do SDD+TDD (Rounds 1--7)}
\label{tab:pre_sdd}
\begin{tabular}{p{0.25\textwidth}p{0.50\textwidth}}
\toprule
\textbf{Característica} & \textbf{Estado Pré-SDD+TDD} \\
\midrule
Metodologia de correção & Reativa: banca simula $\to$ corrige $\to$ repete \\
Garantia contra regressão & Nenhuma — cada correção podia quebrar outras partes \\
Rastreabilidade & Ausente — decisões residiam na memória da sessão \\
Automação & Parcial — dependia de validação humana a cada iteração \\
Confiabilidade & Subjetiva — critério de ``parece bom'' vs ``está provado'' \\
Suítes de teste & 0 — nenhum teste automatizado \\
Score médio & 92,3/100 (Rounds 1--7) \\
\bottomrule
\end{tabular}
\end{table}

\subsection{Round 8: A Introdução do SDD+TDD}

O Round 8 marcou o ponto de inflexão metodológica do ecossistema. Três
inovações foram introduzidas simultaneamente, cada uma endereçando uma das
limitações identificadas:

\subsubsection{Spec-Driven Development (SDD)}

A primeira inovação foi a exigência de que \textbf{toda implementação fosse
precedida por uma especificação formal}. Para o anteprojeto de doutorado
submetido ao programa de pós-graduação, foram produzidas 7 especificações
modulares — cada uma definindo explicitamente: (i) entradas esperadas;
(ii) transformações a serem aplicadas; (iii) saídas esperadas; e (iv)
critérios de aceitação objetivos.

Esta prática, derivada da engenharia de software formal, eliminou a
ambiguidade que causava os ciclos infinitos de correção dos rounds anteriores.
Quando um revisor da banca simulada apontava um problema, a especificação
permitia determinar se o problema era de \textit{implementação} (o código
não seguia a spec) ou de \textit{design} (a spec estava errada) — uma
distinção impossível sem especificações explícitas.

\subsubsection{Test-Driven Development (TDD)}

A segunda inovação foi a inversão do fluxo de trabalho: \textbf{testes antes
da implementação}. Nove critérios de teste (CTs) foram escritos antes que
uma única linha do anteprojeto fosse modificada. Cada CT verificava uma
propriedade específica e incontroversa: ``a seção de metodologia deve
explicitar o paradigma epistemológico adotado'', ``todas as referências
bibliográficas devem possuir DOI verificável'', ``o resumo deve conter
entre 150 e 500 palavras''.

O ciclo RED $\to$ GREEN $\to$ REFACTOR — executar o teste e vê-lo falhar
(RED), implementar a correção mínima necessária para passar (GREEN), e
então melhorar a implementação sem quebrar os testes existentes (REFACTOR)
— substituiu o ciclo anterior de ``escreve tudo $\to$ revisa tudo $\to$
corrige tudo''. A diferença fundamental: no ciclo TDD, cada correção é
\textit{validada automaticamente} antes de ser aceita.

\subsubsection{DecisionNode: Memória Estruturada de Decisões}

A terceira inovação foi o registro formal de decisões arquiteturais via
\textbf{Architecture Decision Records} (ADRs). Três ADRs foram registradas
no Round 8, documentando: (i) a escolha do paradigma SDD+TDD como metodologia
padrão; (ii) a estratégia de teste com 3 quality gates independentes; e
(iii) o protocolo de anonimato para documentos submetidos a avaliação.

Cada ADR contém: contexto (por que a decisão era necessária), decisão (o que
foi decidido), alternativas consideradas (o que foi rejeitado e por quê), e
consequências (o que muda com esta decisão). Esta estrutura garante que
decisões tomadas em um round possam ser compreendidas, reutilizadas e
questionadas em rounds subsequentes.

\subsection{Round 9: A Consolidação em Pipeline Autônomo}

O Round 9 representou a maturação do SDD+TDD de metodologia pontual para
\textbf{pipeline autônomo}. O sistema AutoEvolve — SENSE $\to$ DIAGNOSE
$\to$ FIX $\to$ VERIFY $\to$ EVOLVE $\to$ LEARN — é a materialização do
ciclo TDD em escala de ecossistema.

A Tabela~\ref{tab:tdd_pipeline} mapeia cada fase do TDD clássico para sua
contraparte no pipeline AutoEvolve:

\begin{table}[H]\centering
\caption{Correspondência TDD clássico $\leftrightarrow$ Pipeline AutoEvolve}
\label{tab:tdd_pipeline}
\begin{tabular}{p{0.20\textwidth}p{0.25\textwidth}p{0.35\textwidth}}
\toprule
\textbf{TDD Clássico} & \textbf{AutoEvolve} & \textbf{Função no Ecossistema} \\
\midrule
RED (escrever teste que falha) & DIAGNOSE & Parsear arquivo \texttt{.log} com regex, detectar overfull/underfull boxes, erros LaTeX, font warnings \\
GREEN (implementar mínimo) & FIX & Backup automático + correção seletiva: encurtamento textual, \texttt{\textbackslash sloppy} wrapper, \texttt{\textbackslash raggedright} em colunas \\
REFACTOR (melhorar sem quebrar) & VERIFY & 3 quality gates (Compilation 5, Structure 6, Quality 5) = 16 testes. Se FAIL, re-entra no loop \\
— & EVOLVE & Registrar padrão de correção em \texttt{fix\_history.json}, atualizar catálogo F01--F10 \\
— & LEARN & Gerar insight: tendências, padrões mais frequentes, recomendações \\
\bottomrule
\end{tabular}
\end{table}

O resultado imediato do Round 9 foi a eliminação de 4 overfull boxes (máximo
11,7pt) e 1 underfull box (badness 10000) em \textbf{uma única iteração} do
pipeline. As 3 suítes TDD — Compilation (5/5), Structure (6/6) e Quality
(5/5) — estabeleceram o padrão 16/16 GREEN como quality gate. O catálogo de
10 padrões de correção (F01--F10) documentou estratégias reutilizáveis:
encurtamento textual (F04), \texttt{\textbackslash sloppy} wrapper (F01),
\texttt{\textbackslash raggedright} em colunas (F02), entre outros.

\subsection{O Impacto no CORA-Eval (Rounds 11--14)}

Quando o CORA-Eval foi concebido no Round 11, o SDD+TDD já estava integrado
ao DNA do ecossistema. Esta herança metodológica explica por que o CORA-Eval
não é apenas um catálogo de 150 problemas — é um \textbf{sistema de
verificação cientométrica com rastreabilidade total}.

\subsubsection{Suítes TDD como Evidência}

Cada uma das 10 dimensões do CORA-Eval possui uma suíte de teste automatizada
que segue rigorosamente o ciclo RED $\to$ GREEN $\to$ REFACTOR. Nenhum score
é registrado no \texttt{cora\_scores.json} sem que o teste correspondente
passe. Esta prática — radicalmente diferente de benchmarks como MATH ou MMLU,
onde a ``verificação'' consiste em comparar a resposta final contra um
gabarito — garante que cada afirmação sobre a maturidade científica do
ecossistema é respaldada por evidência executável e reproduzível.

\subsubsection{Pipeline de 5 Estágios}

O pipeline SENSE $\to$ DIAGNOSE $\to$ FIX $\to$ VERIFY $\to$ EVOLVE $\to$
LEARN do Round 9 foi adaptado para o CORA-Eval como Extração $\to$ Resolução
$\to$ Verificação $\to$ Pontuação $\to$ Aprendizado. A estrutura é
isomórfica — apenas o domínio de aplicação mudou de documentos LaTeX para
raciocínio científico —, demonstrando a generalidade do framework SDD+TDD.

\subsubsection{Rastreabilidade Total}

Cada score no \texttt{cora\_scores.json} é rastreável a um teste específico,
em uma suíte TDD específica, com ground truth documentado e verificadores
Cora aplicados. O Anexo D desta dissertação permite a qualquer leitor
recalcular manualmente o CORA-Score e verificar cada contribuição —
propriedade impossível em benchmarks que não adotam TDD.

\subsection{Antes e Depois: A Transformação Quantificada}

A Tabela~\ref{tab:before_after} quantifica o impacto do SDD+TDD no
ecossistema:

\begin{table}[H]\centering
\caption{Impacto do SDD+TDD: antes vs depois}
\label{tab:before_after}
\begin{tabular}{p{0.30\textwidth}p{0.20\textwidth}p{0.20\textwidth}}
\toprule
\textbf{Métrica} & \textbf{Pré-SDD+TDD (R1--R7)} & \textbf{Pós-SDD+TDD (R8--R14)} \\
\midrule
Regressões por correção & $\sim$30\% (estimado) & 0\% (16/16 gates) \\
Tempo médio de correção & 3--5 iterações & 1 iteração \\
Padrões de correção documentados & 0 & 10 (F01--F10) \\
Decisões arquiteturais rastreáveis & 0 & 3 ADRs + 8 snapshots \\
Suítes de teste automatizadas & 0 & 13 (3 LaTeX + 10 CORA) \\
Cobertura de testes & 0\% & 99,0\% (113/114 testes) \\
CORA-Score (maturidade científica) & N/A (não existia) & 2,99 (Pesquisa) \\
Confiança nos resultados & Subjetiva & Objetiva + auditoria externa \\
\bottomrule
\end{tabular}
\end{table}

\subsection{Implicações para a Ciência da Computação}

A experiência do OpenCode Ecosystem com SDD+TDD oferece lições que transcendem
o contexto específico deste projeto. A principal delas é que \textbf{a
metodologia de teste não é um acessório da inteligência artificial — é um
componente fundamental de sua confiabilidade}.

O salto qualitativo entre os Rounds 7 e 14 não foi impulsionado por modelos
maiores, mais dados de treinamento ou arquiteturas mais complexas. Foi
impulsionado pela adoção de uma prática de engenharia de software estabelecida
há décadas — escrever testes antes do código, verificar antes de aceitar,
documentar antes de esquecer — e sua adaptação sistemática ao domínio do
raciocínio científico.

Esta observação tem implicações diretas para o design de sistemas de IA
científica. Se o objetivo é produzir raciocínio confiável e verificável — e
não apenas respostas plausíveis —, então \textbf{o TDD não é opcional. É
infraestrutura}. O CORA-Eval operacionaliza este princípio, fornecendo a
primeira métrica quantitativa do impacto do TDD na maturidade científica de
sistemas de IA.
